% !TeX encoding = UTF-8
% !TeX spellcheck = de_DE

\chapter{Die Axiomatische Grundlegung}
\label{chap:axiome}

\section{Einleitung und Ziel der Theorie}
Dieses Dokument legt die formale Grundlage einer neuartigen Gesellschaftstheorie. Ihr Ziel ist es, die notwendigen und hinreichenden Bedingungen für eine \emph{wirklich demokratische} Gesellschaftsordnung zu definieren. Anders als deskriptive oder historische Theorien geht sie von einem streng normativen und axiomatischen Ansatz aus. Die folgenden fünf Axiome stellen die unverrückbaren Grundpostulate dar, aus denen sich die gesamte Theorie logisch ableitet.

\section{Das Axiomensystem $\mathcal{S}$}
Im Folgenden sei $G$ ein Modell einer Gesellschaft, bestehend aus einer Menge von Individuen $I$, einer Struktur der Entscheidungsfindung und einer Verteilung von Ressourcen und Rechten.

\subsection{Axiom 1: Der normative Kern -- Wirkliche Demokratie}
\label{ax:1}

\textbf{Formale Fassung:}
Sei $\mathcal{D}(G) \in [0,1]$ ein Maß für den Grad der verwirklichten Demokratie in $G$, wobei $\mathcal{D}(G)=1$ einen idealen Zustand bezeichnet. Das erste Axiom postuliert:
\[
\boxed{\forall G: G \text{ soll so beschaffen sein, dass } \mathcal{D}(G) = 1.}
\]

\textbf{Interpretation:}
Dies ist das normative Fundament der Theorie. Der Terminus \glqq wirklich\grqq\ weist darauf hin, dass es nicht um die Erfüllung formaler Kriterien (wie regelmäßige Wahlen) allein geht, sondern um die substanzielle Verwirklichung von \emph{Volkssouveränität}. Dies impliziert die maximale Verwirklichung politischer Gleichheit, gleicher Einflussmöglichkeiten und der Abwesenheit von struktureller Entfremdung der Entscheidungsmacht von den betroffenen Individuen.

\subsection{Axiom 2: Die strukturelle Implikation -- Anti-Parteien-Prinzip}
\label{ax:2}

\textbf{Formale Fassung:}
Sei $P(G)$ die Menge stabiler Organisationen in $G$, deren primärer Zweck die dauerhafte Bündelung und Ausübung politischer Macht ist (im Folgenden \glqq Parteien\grqq). Das zweite Axiom stellt fest:
\[
\boxed{\mathcal{D}(G) = 1 \quad \Rightarrow \quad P(G) = \varnothing.}
\]
Äquivalent: Echte Demokratie verlangt individuelle, unmittelbare Mitbestimmung. Die Existenz von Parteien als Machtkonzentrationen ist undemokratisch.

\textbf{Interpretation:}
Dieses Axiom identifiziert politische Parteien nicht als Vehikel, sondern als Hindernisse der Demokratie. Ihre inhärente Logik der Machtakkumulation, Programmbündelung und Elitenbildung steht im Widerspruch zu individueller Autonomie und unverzerrter Willensbildung. Eine wirklich demokratische Gesellschaft muss daher auf nicht-hierarchischen, direkt-partizipativen oder strikt delegierenden Entscheidungsstrukturen basieren.

\subsection{Axiom 3: Die materielle Bedingung -- Begrenzte Ungleichheit}
\label{ax:3}

\textbf{Formale Fassung:}
Sei $V_i$ das Vermögen (oder die ökonomische Potenz) des Individuums $i \in I$ und $\varepsilon \in \mathbb{R}^+$ eine gesellschaftlich festgelegte, absolute Obergrenze. Das dritte Axiom fordert die Existenz einer Verfassung mit:
\[
\boxed{\forall i \in I: V_i \leq \varepsilon \quad \text{(Begrenzung)} \qquad \text{und} \qquad \forall i,j \in I: R_i = R_j \quad \text{(gleiche Rechte)}.}
\]
Die Parole \glqq Wohlstand ja, Reichtum nein\grqq\ operationalisiert dies: $\varepsilon$ markiert die scharfe Trennlinie zwischen legitimen Wohlstand und illegitimem, demokratiegefährdendem Reichtum.

\textbf{Interpretation:}
Politische Gleichheit (Axiom~\ref{ax:1}) ist ohne materielle Gleichheit nicht denkbar, da ökonomische Macht stets in politischen Einfluss konvertierbar ist. Dieses Axiom schafft die notwendige Voraussetzung für Axiom~\ref{ax:1} und \ref{ax:2}, indem es die Entstehung ökonomischer Eliten und damit neuer politischer Machtasymmetrien strukturell unterbindet.

\subsection{Axiom 4: Das Verfahren -- Epistokratisch-Demokratische Dualität}
\label{ax:4}

\textbf{Formale Fassung:}
Sei $E$ ein evidenzbasierter, wissenschaftlicher Maßnahmenkatalog zu einem politischen Problem. Sei $Z(G)$ der Prozess der kollektiven Zustimmung in $G$. Das vierte Axiom verlangt:
\[
\boxed{\text{Jeder politische Entscheidungsprozess hat die Form: } E \longrightarrow Z(G).}
\]
Weiter gilt: Die Fähigkeit von $G$, $E$ korrekt zu verstehen und zu bewerten, erfordert, dass $G$ eine \emph{Wissensgesellschaft} ist.

\textbf{Interpretation:}
Entscheidungen sollen rational und auf der besten verfügbaren Evidenz basieren. Die finale Autorität liegt jedoch ausschließlich bei der betroffenen Gemeinschaft. Diese Dualität bewahrt vor Technokratie (da die Gemeinschaft zustimmen muss) und vor Irrationalismus (da Vorschläge evidenzbasiert sein müssen). Die Forderung nach einer Wissensgesellschaft ist die funktionale Konsequenz, um informierte Zustimmung zu ermöglichen.

\subsection{Axiom 5: Die Meta-Ebene -- Die Unveränderliche Verfassung}
\label{ax:5}

\textbf{Formale Fassung:}
Sei $C$ die Verfassung der Gesellschaft $G$, welche die Axiome 1-4 operationalisiert. Das fünfte Axiom deklariert:
\[
\boxed{C \text{ ist einfach und unveränderlich.}}
\]
\glqq Einfach\grqq\ bezieht sich auf ihre Klarheit und Beschränkung auf fundamentale Prinzipien. \glqq Unveränderlich\grqq\ bedeutet, dass es keinen legalen Mechanismus zu ihrer Abschaffung oder zur Änderung ihrer axiomatischen Grundsätze gibt.

\textbf{Interpretation:}
Dieses Axiom schützt das System vor seiner eigenen demokratischen Aufhebung. Es friert den revolutionären, demokratischen Moment ein und verhindert einen schleichenden oder gewaltsamen Rückfall in Oligarchie, Autokratie oder unbegrenzten Kapitalismus. Es ist die Garantie für die permanente Gültigkeit der vorherigen Axiome.

\section{Synthese: Das axiomatische System $\mathcal{S}$}
\label{sec:synthese}

Die fünf Axiome bilden ein konsistentes, hierarchisches System:
\begin{itemize}
    \item \textbf{Axiom 1} ist das \emph{normative Ziel}.
    \item \textbf{Axiom 5} ist die \emph{meta-normative Garantie} für dieses Ziel.
    \item \textbf{Axiom 2} und \textbf{3} sind die \emph{strukturellen und materiellen notwendigen Bedingungen} zur Erreichung des Ziels.
    \item \textbf{Axiom 4} ist das \emph{Verfahren}, das den Betrieb des Systems unter Einhaltung aller anderen Axiome regelt.
\end{itemize}

Formal definieren wir das axiomatische System $\mathcal{S}$ als die Konjunktion aller Axiome:
\[
\boxed{\mathcal{S} := \text{A1} \land \text{A2} \land \text{A3} \land \text{A4} \land \text{A5}.}
\]
Eine Gesellschaft $G$, die das Modell $\mathcal{S}$ instanziiert, nennen wir eine $\mathcal{S}$-Gesellschaft. Alle folgenden theoretischen Ableitungen und Analysen beziehen sich auf dieses System $\mathcal{S}$.
